% This is a modified version of Springer's LNCS template suitable for anonymized MICCAI 2025 main conference submissions. 
% Original file: samplepaper.tex, a sample chapter demonstrating the LLNCS macro package for Springer Computer Science proceedings; Version 2.21 of 2022/01/12

\documentclass[runningheads]{llncs}
%
\usepackage[T1]{fontenc}
% T1 fonts will be used to generate the final print and online PDFs,
% so please use T1 fonts in your manuscript whenever possible.
% Other font encodings may result in incorrect characters.
%
\usepackage{graphicx,verbatim}
\usepackage{outlines}

% Used for displaying a sample figure. If possible, figure files should
% be included in EPS format.
%
% If you use the hyperref package, please uncomment the following two lines
% to display URLs in blue roman font according to Springer's eBook style:
%\usepackage{color}
%\renewcommand\UrlFont{\color{blue}\rmfamily}
%\urlstyle{rm}
%
\begin{document}
%
\title{Lesion-Aware Stroke MRI Processing Pipeline using Diffusion Models}
%\titlerunning{Abbreviated paper title}
% If the paper title is too long for the running head, you can set
% an abbreviated paper title here
%
\begin{comment}  %% Removed for anonymized MICCAI submission
\author{First Author\inst{1}\orcidID{0000-1111-2222-3333} \and
Second Author\inst{2,3}\orcidID{1111-2222-3333-4444} \and
Third Author\inst{3}\orcidID{2222--3333-4444-5555}}
%
\authorrunning{F. Author et al.}
% First names are abbreviated in the running head.
% If there are more than two authors, 'et al.' is used.
%
\institute{Princeton University, Princeton NJ 08544, USA \and
Springer Heidelberg, Tiergartenstr. 17, 69121 Heidelberg, Germany
\email{lncs@springer.com}\\
\url{http://www.springer.com/gp/computer-science/lncs} \and
ABC Institute, Rupert-Karls-University Heidelberg, Heidelberg, Germany\\
\email{\{abc,lncs\}@uni-heidelberg.de}}

\end{comment}

\author{Anonymized Authors}  %% Added for anonymized MICCAI submission
\authorrunning{Anonymized Author et al.}
\institute{Anonymized Affiliations \\
    \email{email@anonymized.com}}
  
\maketitle              % typeset the header of the contribution
%
\begin{abstract}
We previously introduced a diffusion-model framework for analyzing and reversing brain lesion effects, which we have adapted here to generate normal-appearing MRIs from stroke MRIs. 
Focal stroke lesions violate the normal-anatomy assumptions
of conventional MRI processing and can corrupt measurements well beyond
the lesion. We evaluate a lesion-aware workflow that automatically
delineates the lesion, replaces a dilated target using a 3D diffusion model,
and applies BrainSuite. Generated tissue is treated only as a computational
scaffold. No intensity or morphometric measurement from any region
overlapping with the inpainting target is used as a biological endpoint.
The utility of the inpainting based workflow is demonstrated by application to aphasia stroke dataset. 
In 211 participants, each T1 acquisition was processed directly and
with the lesion-aware workflow. Among homologous ROI pairs entirely
outside the dilated target, the workflow reduced interhemispheric asymmetry
in tissue volume from 22.3\% to 20.6\% (adjusted p 0.029),
cortical area from 21.4% to 19.0% (adjusted p 0.008), and cortical thickness
from 12.8\% to 10.5% (adjusted p < 0.001).
We also computed the correlation between Left perisylvian lesion burden and Western Aphasia Battery Aphasia Quotient (WAQ) after
adjustment for total lesion volume, age at stroke, and assessment delay. 
Left perisylvian lesion burden was
associated with lower WAQ 
(partial correlation = -0.20, p=0.006). These results support inpainting
as a means to stabilize downstream measurements in spared brain, not
as a method for measuring the invented tissue itself.

This framework treats lesion analysis as an inverse problem by estimating how an observed lesioned brain could have arisen from pre-lesion healthy anatomy. The key conceptual advance is to distinguish irreversibly damaged tissue from surrounding tissue that has been displaced or compressed by lesion growth or swelling, rather than labeling both as a single abnormal region.
The inpainting model was implemented with a U-Net-based 3D diffusion architecture and trained using irregular lesion-like masks generated from perturbed 3D shapes, enabling the model to synthesize realistic tissue fills guided by surrounding anatomy. The trained model replaces masked lesion voxels with plausible healthy tissue consistent with surrounding anatomy. Registration between the original and inpainted images yields a local deformation field, allowing the inverse field to estimate tissue displacement near the lesion while preserving unchanged anatomy elsewhere.

\keywords{stroke  \and inpainting \and diffusion models \and MRI}
% Authors must provide keywords and are not allowed to remove this Keyword section.

\end{abstract}
\section{Introduction}
\begin{outline}[enumerate] % Uses numbered/alphanumeric formatting
    \1 Literature survey
        \2 freesurfer paper
            \3 how did they validate?
            \3 limitations?
        \2 other tools? other in-painting?
    \1 Relevance of problem
\end{outline}
We introduced a diffusion-model framework for analyzing and reversing brain lesion effects \cite{Anonymous,Jyothi2023}, which we have adapted here to generate normal-appearing MRIs from stroke MRIs. This framework treats lesion analysis as an inverse problem by estimating how an observed lesioned brain could have arisen from pre-lesion healthy anatomy. The key conceptual advance is to distinguish irreversibly damaged tissue from surrounding tissue that has been displaced or compressed by lesion growth or swelling, rather than labeling both as a single abnormal region.
The inpainting model was implemented with a U-Net-based 3D diffusion architecture and trained using irregular lesion-like masks generated from perturbed 3D shapes, enabling the model to synthesize realistic tissue fills guided by surrounding anatomy. The trained model replaces masked lesion voxels with plausible healthy tissue consistent with surrounding anatomy. Registration between the original and inpainted images yields a local deformation field, allowing the inverse field to estimate tissue displacement near the lesion while preserving unchanged anatomy elsewhere (see Fig. \ref{fig2}). Another diffusion-based inpainting approach (FastSurfer-LIT) was recently developed by \cite{Pollak2025}, with results demonstrating the value of diffusion-based lesion inpainting for restoring segmentation and surface reconstruction in abnormal brains. Our approach differs from the one in FastSurfer-LIT because we model tissue deformation as part of the inpainting process and propagate uncertainty information into downstream processing, extending foundational diffusion models for medical anomaly detection \cite{Wyatt2022}. 

\section{Methods}
\begin{figure}[t!]
\includegraphics[width=\textwidth]{figs/Fig1_24pt.png}
\caption{{\bf Inpainting of lesions.} (A) original T1, (B) automatically delineated stroke lesion; (C) uncertainty map; (D) deformation corrected and inpainted T1; (E) pressure estimate from divergence of the Jacobian of the deformation field; (F) MRI volumetrically labeled using BrainSuite; (G) surface labels; (H) cortical thickness map.} \label{fig2}
\end{figure}

\subsection{Geometric Modeling of Brain Lesion Effects}
We model the lesion process as a combination of irreversible tissue damage (the core lesion) and surrounding tissue deformation (mass effect). To reverse this, we first identify the lesion area and then estimate the deformation field that restores the displaced tissue to its original configuration. This is achieved by registering the lesioned brain image to an inpainted version of the image where the lesion is replaced by healthy-appearing tissue. The resulting deformation field isolates the irreversible damage to a core region while correcting the mass effect on the surrounding anatomy  \cite{Vercauteren2009}.


\subsection{Diffusion-based Inpainting}
A core component of our pipeline is a 3D inpainting diffusion model, which is used to replace lesion abnormalities with healthy-appearing tissue. We developed a guided inpainting module based on a 3D U-Net diffusion architecture, building upon recent advances in generative models for brain imaging \cite{Pinaya2022}. The model is fine-tuned to accept an encoding mask (the lesion mask) alongside the noisy image. During training, we generate irregular 3D lesion-like masks using spherical meshes perturbed by Perlin noise to simulate natural lesion variability. The model learns to synthesize plausible healthy tissue in the masked regions conditioned on the surrounding unmasked anatomy.


\subsection{Uncertainty Propagation}
As diffusion models are generative and probabilistic, multiple inpainting samples can be generated for a single lesion. This allows us to compute voxel-wise variance across the generated samples, producing an uncertainty map (as shown in Fig. \ref{fig2}C). This uncertainty information can be propagated into downstream processing steps, enabling probabilistic segmentation and more robust anatomical labeling that accounts for the ambiguity of the imputed tissue.

\subsection{Integration into Conventional Pipeline}
To enable standard neuroimaging software to process stroke MRIs, we integrate our lesion reversal pipeline as a preprocessing step. The stroke MRI is passed through the inpainting model to generate a healthy-appearing brain, which is then processed using standard anatomical pipelines (e.g., BrainSuite \cite{Shattuck2002}) for skull stripping, tissue classification, and surface extraction, mitigating failures common to automated stroke segmentation tools \cite{Pustina2016}. The resulting labels and surfaces can then be mapped back into the original lesioned brain space using the inverse of the estimated deformation field, resolving classical issues associated with spatial normalization in the presence of focal lesions \cite{Brett2001}.

\subsection{Evaluation}

\subsubsection{ARC Data and Processing}
We used the Aphasia Recovery Cohort (ARC; OpenNeuro ds004884, version 1.0.2) for this study \cite{gibson2024aphasia}. ARC is a longitudinal, BIDS-formatted repository comprising demographic, behavioral, and multimodal MRI data from 230 individuals with chronic left-hemisphere stroke, collected for 229 participants across 902 imaging sessions. MRI were acquired on Siemens 3-T systems and included T1, T2, FLAIR, DWI and task fMRI images. The dataset also provides expert-defined lesion masks based on first-visit T2-weighted images. 

For this study, we did not rely on the provided manual delineation, but used our trained nn-UNet to delineate the lesions. The lesion boundary was then dilated using 3mm spherical stencil and the corresponding region was inpainted using our inpainting method.
Both the original and inpainted data were separately processed using BrainSuite Anatomical Pipeline. The pipeline includes skull extraction, bias field correction, tissue classification, cerebrum/cerebellum identification, inner and pial -cortical surface extraction. This is followed by Surface Constrained Volumetric Registration sequence that aligns the brain with the BrainSuiteAtlas1 anatomical atlas. 

We analyzed the intersection of completed original and inpainted BrainSuite outputs in the ARC BIDS archive: [TBD] unique participants and one T1 acquisition per participant. Median automatically delineated lesion volume was [TBD]~mL (IQR [TBD]--[TBD]); [TBD] lesions were left-dominant and [TBD] was right-dominant. Participant-level WAB Aphasia Quotient, age at stroke, and days from stroke to WAB were obtained from the participant table. No pre-stroke T1 or expert BrainSuite labels were available, so all endpoints use internal controls and are interpreted as plausibility or construct-validity measures rather than reconstruction accuracy.





\subsubsection{Spared-Anatomy Endpoints}
The generated region was excluded from all morphometric endpoints. BrainSuite volumetric labels \(1000+\mathrm{ROI}\) and \(2000+\mathrm{ROI}\) (cortical gray and white matter) were first reduced to their canonical ROI identifiers. A homologous ROI pair was eligible only when neither hemisphere had any voxel overlap with the 3-mm-dilated inpainting target; even a single overlapping voxel excluded the pair. Thus, synthetic tissue and the blending margin contributed no regional volume, area, or thickness value to the analysis.

For each eligible spared pair, interhemispheric asymmetry was \(A=200|x_L-x_R|/(x_L+x_R)\) for total tissue volume, mid-cortical area, and mean cortical thickness. The median across eligible pairs formed one participant-level observation for each endpoint. This symmetry measure does not assume that the generated tissue is correct; it tests whether using inpainting as a preprocessing scaffold produces more internally consistent measurements elsewhere in the brain.

\subsubsection{Aphasia Construct Validity}
For a separate construct-validity analysis, we intersected the original binary source-lesion mask with post-inpainting atlas labels in a prespecified left perisylvian network: pars opercularis, pars triangularis, pars orbitalis, supramarginal and angular gyri, superior, transverse, and middle temporal gyri, and insula. This analysis uses lesion presence and anatomical localization only; it does not use generated intensity, thickness, area, or volume as a biological measurement. Partial Spearman correlation related localized lesion burden to WAB Aphasia Quotient while controlling for total lesion volume, age at stroke, and \(\log(1+\mathrm{days})\). Total lesion volume was evaluated while controlling for age and assessment delay; the right homolog network was a negative control. This analysis establishes clinical plausibility of the anatomical localization and is not treated as a paired superiority endpoint.

Paired pipeline endpoints were tested with two-sided Wilcoxon signed-rank tests. We report the paired median improvement, a participant-bootstrap 95\% confidence interval (10,000 resamples), and rank-biserial effect size (positive favors inpainting). Holm correction was applied across the three spared-anatomy endpoints and separately across three clinical associations. Analyses were performed with Python, NumPy, SciPy, pandas, and NiBabel; the complete script and case-level tables accompany the manuscript.



\subsubsection{Evaluation Data}
The inpainting diffusion model was trained on T1-weighted images from healthy cohorts. For evaluation on stroke MRIs, we utilize datasets of stroke patients with manually delineated lesion masks (e.g., ATLAS dataset \cite{Liew2022} and ISLES benchmarks \cite{Maier2015}) to assess the pipeline's ability to restore pre-lesion anatomy and improve automated labeling accuracy compared to direct processing of the lesioned brains.





\subsubsection{Evaluation Criteria}
We evaluate the performance of our pipeline using both geometric and anatomical metrics. For inpainting quality, we calculate the Structural Similarity Index (SSIM) and Normalized Mean Squared Error (NMSE) between the inpainted regions and ground truth (using synthetically masked healthy brains). To assess the improvement in downstream processing, we compare the accuracy of anatomical labels generated by our pipeline against expert manual labels, measuring the Dice similarity coefficient in regions adjacent to the stroke lesion.

\subsection{Aphasia quotient prediction}
The original binary lesion, never the generated tissue, was intersected with
post-inpainting atlas labels. A prespecified left perisylvian network contained
pars opercularis, triangularis, and orbitalis; supramarginal and angular gyri;
superior, transverse, and middle temporal gyri; and insula. Partial Spearman
correlation related left-network lesion burden to AQ while controlling for
total lesion volume, age at stroke, and $\log(1+\mathrm{WAB\ days})$. Total
volume controlled for age and assessment delay; the right homolog network was
a negative anatomical control.

\subsection{Lesion-associated deformation }
SVReg's inverse coordinate map $F(\vect{x})$ gives the subject-voxel coordinate
corresponding to atlas voxel $\vect{x}$. On valid contralateral brain voxels, a
robust affine map $F_A(\vect{x})=M\vect{x}+\vect{b}$ is fitted with iterative
median-absolute-deviation outlier rejection. Removing this global component and
mapping its residual back to atlas axes gives
\begin{equation}
 \vect{u}(\vect{x})=M^{-1}\{F(\vect{x})-F_A(\vect{x})\},
 \label{eq:residual}
\end{equation}
expressed in atlas-aligned millimeters and smoothed with a 2-mm Gaussian kernel.
Let $R$ reflect position through the atlas midline and reflect a vector's
left--right component. The lesion-associated asymmetry field is
\begin{equation}
 \vect{u}_{L}(\vect{x})=\vect{u}(\vect{x})-R\vect{u}(R\vect{x}).
 \label{eq:mirror}
\end{equation}
Both members of a mirrored pair must lie in the atlas and mapped subject brain
masks. The lesion, dilated target, and reflected target are excluded. Cases
with lesion laterality below 0.80 are flagged as unsupported.

Magnitude is $m=\|\vect{u}_L\|$. Outside the lesion, the gradient of its signed
distance map gives outward unit normal $\vect{n}$; radial displacement is
$r=\vect{u}_L\cdot\vect{n}$, with $r>0$ outward and $r<0$ inward. Local volume
change uses $J=\det(I+\nabla\vect{u})$ and
$\logJ(\vect{x})=\log J(\vect{x})-\log J(R\vect{x})$, following the
deformation-based morphometry interpretation of Jacobians
\cite{Chung2001}. Nonpositive Jacobians are invalid. Magnitude, radial, and
$\logJ$ are summarized in 3--5, 5--10, 10--20, and 20--40~mm shells. Although
Jacobian deformation has been related to radiologist-rated tumor mass effect
\cite{Zacharaki2009}, a chronic infarct is not space occupying; we therefore
avoid a mechanistic pressure claim.

The direct, non-inpainted SVReg branch is normalized identically. The magnitude
$\|\vect{u}_{L,\mathrm{inp}}-\vect{u}_{L,\mathrm{raw}}\|$ is stored as
registration sensitivity, not biological deformation. This comparison and
the exclusion of synthetic voxels address the known tendency of lesions to
bias nonlinear registration \cite{Brett2001,Sdika2009}.

\subsection{Study 4: incremental AQ prediction}
Table~\ref{tab:models} defines prespecified ridge-regression models. The primary
contrast asks whether deformation improves prediction beyond age, assessment
delay, total lesion volume, laterality, and left/right language-network lesion
burden. Eight compact 3--20~mm features are used: median and 95th-percentile
magnitude; median and mean-absolute radial displacement; outward and inward
radial integrals; and positive and negative log-Jacobian integrals. A secondary
model adds raw--inpainted sensitivity and contralateral affine-fit RMSE to test
whether apparent benefit is registration-driven. If uncertainty coverage is at
least 90\%, soft volume, entropy, boundary ambiguity, and threshold-sensitivity
features are added to the conventional lesion model, followed by deformation.

\begin{table}[t]
\caption{Prespecified AQ prediction models, evaluated on one identical
QC-passing cohort. Features are cumulative from top to bottom within each
comparison.}
\label{tab:models}
\centering
\scriptsize
\begin{tabularx}{\textwidth}{p{0.28\textwidth}X}
\toprule
Model & Predictors \\
\midrule
Clinical & Age at stroke; $\log(1+\mathrm{WAB\ days})$ \\
Conventional lesion & Clinical; lesion volume; left/right language-network burden; laterality \\
Lesion + deformation & Conventional lesion; magnitude, radial, and log-Jacobian summaries \\
+ registration QC & Above; raw--inpainted field sensitivity; affine-fit RMSE \\
Lesion + uncertainty & Conventional lesion; expected volumes, entropy mass, boundary/candidate entropy, threshold volumes \\
Uncertainty + deformation & Lesion + uncertainty; deformation summaries \\
\bottomrule
\end{tabularx}
\end{table}

Eligible fields require supported laterality, folding fraction $\leq0.05$, and
at least 1,000 valid voxels 3--20~mm from the lesion. Missing predictors are
median-imputed, standardized, and fitted with ridge regression entirely inside
each training fold. Five-fold outer cross-validation is repeated 20 times; a
four-fold inner loop selects $\alpha\in10^{-4},\ldots,10^4$ by mean absolute
error (MAE). Nested evaluation avoids parameter-selection bias
\cite{Varma2006}. Outcomes are out-of-fold MAE, RMSE, $R^2$, and Pearson $r$.
The primary paired advantage is each participant's mean absolute error under
the conventional model minus that under the deformation model; positive values
favor deformation. We prespecify participant bootstrap confidence intervals,
two-sided paired Wilcoxon tests, and Holm correction \cite{Holm1979}.

\subsection{Statistical analysis for completed endpoints}
Study~1 used two-sided paired Wilcoxon tests, participant-bootstrap 95\%
confidence intervals (10,000 resamples), and rank-biserial effect sizes.
Study~2 used bootstrap intervals for partial Spearman correlations. Holm
correction was applied separately across the three morphometric endpoints and
the three clinical associations. All scripts preserve case-level inputs,
out-of-fold predictions, exact configurations, and subject QC images.





\section{Results}

% \begin{figure}
% \includegraphics[height=0.9\textheight]{figs/Aim1Pipeline_v4.pdf}
% \caption{Pipeline.} \label{fig1}
% \end{figure}

Applying our pipeline to stroke MRIs demonstrated significant improvements in both visual quality and quantitative metrics. The inpainting diffusion model successfully generated realistic, coherent brain tissue that seamlessly blended with the surrounding anatomy, achieving high SSIM and low NMSE on synthetic evaluations. When integrated into the BrainSuite pipeline, our method successfully corrected mislabeling errors commonly observed when conventional tools are applied directly to lesioned brains. Specifically, cortical surface extraction and volumetric labeling in regions adjacent to the stroke lesions were substantially more accurate, as the initial processing was performed on the healthy-appearing estimated brain before being mapped back to the subject's native space.

Representative cases at the 10th, 50th, and 90th percentiles of lesion volume illustrate the downstream processing behavior (Fig.~\ref{fig:representative-outputs}). With increasing lesion burden, direct processing shows progressively greater distortions of volumetric and midcortical surface extraction and labeling, whereas those extracted using inpainted workflow is relatively unaffected by the lesion burdon. The labels within inpainted volume are qualitative pipeline illustrations only. Those labels or the metrics extracted within the outlined generated region are not interpreted as recovered anatomy and do not enter any endpoint analysis.


The study-relevant endpoints all favored the inpainted workflow (Fig.~\ref{fig:arc-validation}A--C; Table~\ref{tab:paired-results}). In spared ROI pairs with zero target overlap, median tissue-volume asymmetry decreased from \ArcSparedVolumeBeforeMedian\% to \ArcSparedVolumeAfterMedian\% (paired reduction \ArcSparedVolumeDeltaMedian{} points, 95\% CI \ArcSparedVolumeDeltaCILow--\ArcSparedVolumeDeltaCIHigh; rank-biserial \(r=\ArcSparedVolumeEffect\), adjusted \(p\) \ArcSparedVolumeP). Cortical-area asymmetry decreased from \ArcSparedAreaBeforeMedian\% to \ArcSparedAreaAfterMedian\% (reduction \ArcSparedAreaDeltaMedian{} points, 95\% CI \ArcSparedAreaDeltaCILow--\ArcSparedAreaDeltaCIHigh; \(r=\ArcSparedAreaEffect\), adjusted \(p\) \ArcSparedAreaP). The largest improvement was in cortical-thickness asymmetry, which decreased from \ArcSparedThicknessBeforeMedian\% to \ArcSparedThicknessAfterMedian\% (reduction \ArcSparedThicknessDeltaMedian{} points, 95\% CI \ArcSparedThicknessDeltaCILow--\ArcSparedThicknessDeltaCIHigh; \(r=\ArcSparedThicknessEffect\), adjusted \(p\) \ArcSparedThicknessP). None of these endpoints includes an ROI touching the inpainting target.

Total lesion volume was strongly associated with lower WAB score after adjustment for age and assessment delay (partial \(\rho=\ArcClinicalVolumeRho\), 95\% CI [TBD] to [TBD]; adjusted \(p\) [TBD]). Using only the original lesion mask and atlas location, left perisylvian lesion burden retained an association after additional adjustment for total lesion volume (partial $\(\rho=[TBD])$, 95\% CI \ArcClinicalLeftCILow{} to [TBD]; adjusted $\(p\)$ \ArcClinicalLeftP; Fig.~\ref{fig:arc-validation}D). The right homolog control was null (\(\rho=\ArcClinicalRightRho\), 95\% CI \ArcClinicalRightCILow{} to \ArcClinicalRightCIHigh; adjusted \(p\) \ArcClinicalRightP), although its interpretation is limited by the almost exclusively left-hemisphere cohort.




\section{Conclusion}
In this work, we presented a lesion-aware processing pipeline that utilizes diffusion models to reverse the effects of stroke lesions on brain MRIs. By separating irreversible tissue damage from reversible deformation and generating healthy-appearing pre-lesion images, our framework enables standard neuroimaging tools such as BrainSuite to reliably process stroke data. This pipeline represents a critical step towards robust automated analysis of pathological brains, facilitating more accurate clinical assessments and large-scale neuroimaging studies of stroke populations.

\begin{credits}
\subsubsection{\ackname} This study was funded by Anonymized (grant number Anonymized).

\subsubsection{\discintname}
The authors have no competing interests to declare that are
relevant to the content of this article.
\end{credits}

\bibliographystyle{splncs04}
\bibliography{lesionaware}
%
% ---- Bibliography ----
%
% BibTeX users should specify bibliography style 'splncs04'.
% References will then be sorted and formatted in the correct style.
%
% \bibliographystyle{splncs04}
% \bibliography{mybibliography}
%
% \begin{thebibliography}{8}
% \bibitem{ref_article1}
% Author, F.: Article title. Journal \textbf{2}(5), 99--110 (2016)

% \bibitem{ref_lncs1}
% Author, F., Author, S.: Title of a proceedings paper. In: Editor,
% F., Editor, S. (eds.) CONFERENCE 2016, LNCS, vol. 9999, pp. 1--13.
% Springer, Heidelberg (2016). \doi{10.10007/1234567890}

% \bibitem{ref_book1}
% Author, F., Author, S., Author, T.: Book title. 2nd edn. Publisher,
% Location (1999)

% \bibitem{ref_proc1}
% Author, A.-B.: Contribution title. In: 9th International Proceedings
% on Proceedings, pp. 1--2. Publisher, Location (2010)

% \bibitem{ref_url1}
% LNCS Homepage, \url{http://www.springer.com/lncs}, last accessed 2023/10/25
% \end{thebibliography}
\end{document}
